\begin{ExerciseBlock}
\ExerciseHeading{06}{Reibung an schiefer Ebene}{work_and_energy_friction_on_an_incline}
Eine Kugel mit Masse $m = 1,5\,$kg rollt aus einer Höhe $h_0 = 1,2\,$m eine schiefe Ebene mit einem Winkel von $\alpha = 37^\circ$ herunter (siehe Abbildung). Dabei tritt eine Rollreibung $\mu_R = 0,05$ auf. \\
Kurz nach Ende der schiefen Ebene staucht die Kugel eine entspannte Feder. Vernachlässigen Sie im Folgenden die Ausdehnung und eine mögliche Rotation der Kugel.

\begin{center}
	\CachedTikzFigure{content/resources/chapters/work_and_energy/tikz/friction_on_an_incline.tex}
\end{center}

\begin{enumerate}
	\item Bestimmen Sie die Geschwindigkeit $v_F$, die die Kugel am Ende der schiefen Ebene erreicht.

	\item Die Feder wird um $\Delta s = 3,0\,$cm zusammengestaucht. Bestimmen Sie die Federkonstante $D$ der
Feder.

	\item Nachdem die Kugel die Feder zusammengedrückt hat, wird die Kugel zurückgeschossen und rollt die schiefe Ebene wieder hoch. Welche Höhe $h_1$ erreicht die Kugel?

	\item Dieser Vorgang wiederholt sich mehrmals. Berechnen Sie die gesamte Wegstrecke, die die Kugel bei diesem Vorgang zurücklegt.
\end{enumerate}
\end{ExerciseBlock}
\UseExerciseSolution{work_and_energy_friction_on_an_incline}
